% Paquetes requeridos para compilar este fragmento:
%   \usepackage{booktabs}   -- tablas formales con \toprule, \midrule, \bottomrule
%   \usepackage{hyperref}   -- si se desea hacer clic en la referencia al CSV
%   \usepackage{amsmath}    -- para notación matemática en línea

% Este archivo se inserta dentro de un capítulo existente.
% NO incluye \documentclass ni \begin{document}.

\section{Conclusión del experimento SA -- instancias pequeñas}
\label{sec:conclusion_sa_small}

\subsection{Configuración experimental}
\label{subsec:config_experimental}

El Recocido Simulado (SA) evaluado en este experimento opera con parámetros de temperatura
calculados automáticamente por instancia: la temperatura inicial sigue la fórmula
\(T_{\text{init}} = 20 \cdot d_{\max} / n\), donde \(d_{\max}\) es el costo máximo de
servicio y \(n\) el número de tareas requeridas. La cadena de Markov (longitud \(L\)) y
la temperatura mínima se derivan de la dinámica de aceptación, sin requerir calibración
manual por parte del usuario.

Los parámetros fijos del experimento son:

\begin{itemize}
  \item \(\alpha_{\text{inter}} = 0{,}8\): factor de enfriamiento para operadores
        inter-ruta.
  \item \texttt{patience} \(= 10\): niveles de temperatura sin mejora antes de ejecutar
        un reheat.
  \item \texttt{reheat\_factor} \(= 0{,}5\): la temperatura asciende al \(50\,\%\) de
        \(T_{\text{init}}\) en cada reheat.
  \item \texttt{max\_reheats\_sin\_mejora} \(= 5\): criterio de parada temprana; cinco
        reheats consecutivos sin mejorar la mejor solución global detienen la ejecución.
  \item 9 operadores de vecindad: \texttt{relocate\_intra}, \texttt{swap\_intra},
        \texttt{2opt\_intra}, \texttt{relocate\_inter}, \texttt{swap\_inter},
        \texttt{2opt\_star}, \texttt{cross\_exchange}, \texttt{or\_opt\_2} y
        \texttt{or\_opt\_3}.
\end{itemize}

El grid search barrió dos hiperparámetros libres: \(\alpha\) (tasa de enfriamiento
intra-ruta) con 11 valores en \([0{,}80,\, 0{,}99]\), y \(p_{\text{inter}}\)
(probabilidad de seleccionar un operador inter-ruta en cada iteración) con 5 valores
en \([0{,}4,\, 0{,}8]\).

% --------------------------------------------------------------------------
\subsection{Cobertura del grid search}
\label{subsec:cobertura_grid}

El diseño experimental combina 11 valores de \(\alpha\) \(\times\) 5 valores de
\(p_{\text{inter}}\) \(\times\) 23 instancias \(\times\) 2 repeticiones, sumando
\textbf{2\,530 corridas} en total. Cada celda del grid (par
\(\alpha \times p_{\text{inter}}\)) acumula 46 corridas (2 repeticiones \(\times\) 23
instancias), y cada instancia es evaluada bajo 110 combinaciones distintas de
hiperparámetros. Esta cobertura permite estimar el efecto marginal de cada
hiperparámetro con márgenes estadísticos razonables, separando la varianza atribuible
a \(\alpha\), a \(p_{\text{inter}}\) y a la variabilidad entre instancias.

Las 23 instancias provienen de dos familias de referencia del CARP: 17 instancias
pequeñas de la familia \textbf{GDB} (Golden, DeArmon \& Baker) y 6 instancias de la
familia \textbf{KSHS}, todas con BKS (\emph{Best Known Solutions}) publicados, lo que
permite calcular gaps absolutos respecto al óptimo conocido.

% --------------------------------------------------------------------------
\subsection{Desempeño global}
\label{subsec:desempeno_global}

Las 2\,530 corridas arrojaron una \textbf{tasa de factibilidad del \(100\,\%\)}: ninguna
corrida terminó con una solución que violara restricciones de capacidad. Esto confirma
que el mecanismo de construcción y los operadores de vecindad implementados preservan
la factibilidad a lo largo de toda la búsqueda.

Los tiempos de ejecución por corrida muestran una distribución asimétrica: mediana de
\textbf{\(3{,}64\) s}, promedio de \textbf{\(8{,}65\) s}, mínimo de \(0{,}19\) s y
máximo de \(112{,}55\) s. La separación entre mediana y promedio refleja la existencia
de instancias con espacios de búsqueda sustancialmente más grandes que detienen el
algoritmo tardíamente.

El número de enfriamientos ejecutados (descensos completos de temperatura) muestra un
rango de 11 a 348, con promedio de 67. Que la mediana de enfriamientos sea notablemente
inferior al máximo indica que la \textbf{parada temprana por
\texttt{max\_reheats\_sin\_mejora} \(= 5\) actúa con frecuencia}: la mayoría de las
corridas no agotan el presupuesto teórico de temperatura, sino que terminan al detectar
estancamiento. Esto es deseable en términos de eficiencia, pero también señala que el
presupuesto efectivo de búsqueda puede ser insuficiente para escapar de óptimos locales
profundos.

El número de iteraciones totales por corrida (mínimo 1\,331, promedio 33\,447, máximo
378\,972) refleja directamente la heterogeneidad en el tamaño de las instancias y en
cuántos niveles de temperatura se recorren antes de la parada.

% --------------------------------------------------------------------------
\subsection{Análisis del grid search: efecto de los hiperparámetros}
\label{subsec:analisis_grid}

\subsubsection{Efecto de \texorpdfstring{$\alpha$}{alpha} (tasa de enfriamiento intra-ruta)}
\label{subsubsec:efecto_alpha}

El parámetro \(\alpha\) muestra un efecto \textbf{monótono y claro}: a mayor valor
(enfriamiento más lento), menor gap promedio frente al BKS. La
Tabla~\ref{tab:efecto_alpha} resume el gap promedio y mediano para cada nivel de
\(\alpha\), promediando sobre las 230 corridas correspondientes (5 valores de
\(p_{\text{inter}}\) \(\times\) 23 instancias \(\times\) 2 repeticiones).

\begin{table}[htbp]
  \centering
  \caption{%
    Gap promedio y mediano por nivel de \(\alpha\), promediado sobre las
    230~corridas de cada nivel (familias GDB y KSHS combinadas).%
  }
  \label{tab:efecto_alpha}
  \begin{tabular}{r r r}
    \toprule
    \(\alpha\) & Gap promedio & Gap mediano \\
    \midrule
    0{,}80 & \(33{,}98\,\%\) & \(33{,}08\,\%\) \\
    0{,}82 & \(33{,}84\,\%\) & \(32{,}89\,\%\) \\
    0{,}84 & \(33{,}86\,\%\) & \(32{,}79\,\%\) \\
    0{,}86 & \(33{,}45\,\%\) & \(32{,}83\,\%\) \\
    0{,}88 & \(33{,}23\,\%\) & \(32{,}73\,\%\) \\
    0{,}90 & \(32{,}92\,\%\) & \(32{,}28\,\%\) \\
    0{,}92 & \(32{,}56\,\%\) & \(31{,}94\,\%\) \\
    0{,}94 & \(32{,}21\,\%\) & \(31{,}86\,\%\) \\
    0{,}96 & \(31{,}74\,\%\) & \(31{,}55\,\%\) \\
    0{,}98 & \(31{,}20\,\%\) & \(31{,}86\,\%\) \\
    \textbf{0{,}99} & \(\mathbf{30{,}89\,\%}\) & \(\mathbf{29{,}93\,\%}\) \\
    \bottomrule
  \end{tabular}
\end{table}

La diferencia entre el extremo inferior (\(\alpha = 0{,}80\), gap promedio \(33{,}98\,\%\))
y el superior (\(\alpha = 0{,}99\), gap promedio \(30{,}89\,\%\)) es de aproximadamente
\textbf{3{,}1 puntos porcentuales}. Aunque la magnitud no es dramática en términos
absolutos, la tendencia es estadísticamente limpia y sugiere que el SA se beneficia de
explorar el espacio de soluciones de forma más pausada: un enfriamiento lento mantiene
temperaturas altas por más iteraciones, lo que favorece la aceptación de movimientos
deteriorantes y reduce la probabilidad de quedar atrapado en mínimos locales tempranos.

\subsubsection{Efecto de \texorpdfstring{$p_{\text{inter}}$}{p\_inter} (probabilidad inter-ruta)}
\label{subsubsec:efecto_pinter}

A diferencia de \(\alpha\), el parámetro \(p_{\text{inter}}\) muestra un efecto
\textbf{pequeño y no monótono}. La variación total entre los cinco niveles evaluados es
inferior a \textbf{\(0{,}3\) puntos porcentuales} en gap promedio
(Tabla~\ref{tab:efecto_pinter}).

\begin{table}[htbp]
  \centering
  \caption{%
    Gap promedio y mediano por nivel de \(p_{\text{inter}}\), promediado sobre las
    253~corridas de cada nivel (familias GDB y KSHS combinadas, todos los valores
    de \(\alpha\)).%
  }
  \label{tab:efecto_pinter}
  \begin{tabular}{r r r}
    \toprule
    \(p_{\text{inter}}\) & Gap promedio & Gap mediano \\
    \midrule
    0{,}4 & \(32{,}89\,\%\) & \(32{,}73\,\%\) \\
    0{,}5 & \(32{,}70\,\%\) & \(32{,}31\,\%\) \\
    \textbf{0{,}6} & \(\mathbf{32{,}59\,\%}\) & \(\mathbf{31{,}86\,\%}\) \\
    0{,}7 & \(32{,}71\,\%\) & \(31{,}86\,\%\) \\
    0{,}8 & \(32{,}69\,\%\) & \(31{,}86\,\%\) \\
    \bottomrule
  \end{tabular}
\end{table}

El mínimo marginal aparece en \(p_{\text{inter}} = 0{,}6\), pero los niveles
\(0{,}6\), \(0{,}7\) y \(0{,}8\) producen resultados prácticamente idénticos. Esto
indica que, dentro del rango explorado, el equilibrio entre vecindarios intra-ruta e
inter-ruta no es un factor determinante del desempeño, al menos cuando el conjunto de
instancias incluye grafos de tamaño pequeño con pocas rutas por solución.

\subsubsection{Configuraciones destacadas}
\label{subsubsec:configs_destacadas}

Las cinco mejores combinaciones \(\alpha \times p_{\text{inter}}\) por gap promedio
sobre las 46 corridas de cada celda se presentan en la Tabla~\ref{tab:configs_destacadas}.

\begin{table}[htbp]
  \centering
  \caption{%
    Cinco mejores configuraciones del grid search, ordenadas por gap promedio ascendente.
    Cada celda acumula 46~corridas (2~repeticiones \(\times\) 23~instancias).%
  }
  \label{tab:configs_destacadas}
  \begin{tabular}{r r r r r}
    \toprule
    \(\alpha\) & \(p_{\text{inter}}\) & Gap promedio & Gap mediano & Desv. estándar \\
    \midrule
    0{,}99 & 0{,}7 & \(30{,}78\,\%\) & \(28{,}61\,\%\) & 10{,}67 \\
    0{,}99 & 0{,}6 & \(30{,}87\,\%\) & \(30{,}77\,\%\) & 10{,}77 \\
    0{,}99 & 0{,}5 & \(30{,}92\,\%\) & \(31{,}04\,\%\) & 10{,}87 \\
    0{,}99 & 0{,}8 & \(30{,}93\,\%\) & \(29{,}92\,\%\) & 10{,}58 \\
    0{,}99 & 0{,}4 & \(30{,}95\,\%\) & \(29{,}24\,\%\) & 10{,}77 \\
    \bottomrule
  \end{tabular}
\end{table}

Las 11 configuraciones con menor gap promedio corresponden a \(\alpha = 0{,}99\) o
\(\alpha = 0{,}98\). Este resultado corrobora que \textbf{\(\alpha\) domina la varianza
del desempeño del grid}, mientras que \(p_{\text{inter}}\) actúa como factor secundario
cuya elección precisa tiene un impacto menor al décimo de punto porcentual en la media.

% --------------------------------------------------------------------------
\subsection{Resultados por instancia}
\label{subsec:resultados_instancia}

Los gaps individuales de la mejor corrida por instancia (consultables en
\texttt{resumen\_mejores\_corridas.csv}, véase también la
Tabla~\ref{tab:sa_small_resumen}) oscilan entre \textbf{\(14{,}13\,\%\)} y
\textbf{\(52{,}00\,\%\)}, con promedio de \(30{,}54\,\%\) y mediana de
\(27{,}59\,\%\). La dispersión es elevada, lo que refleja diferencias estructurales
importantes entre instancias.

Las instancias con menor gap son kshs2 (\(14{,}13\,\%\)), gdb19 (\(14{,}55\,\%\)),
kshs1 (\(15{,}27\,\%\)), gdb13 (\(19{,}03\,\%\)), gdb16 (\(19{,}69\,\%\)) y gdb17
(\(19{,}78\,\%\)). Las instancias de la familia KSHS y algunas GDB con topología más
densa (gdb13, gdb16, gdb17, gdb19) parecen ofrecer un paisaje de búsqueda donde el SA
encuentra soluciones relativamente cercanas al BKS. En estos casos, la combinación de
operadores intra e inter-ruta logra recombinar las rutas de forma efectiva.

En el extremo opuesto, cinco instancias superan el \(40\,\%\) de gap: gdb3
(\(49{,}09\,\%\)), gdb10 (\(52{,}00\,\%\)), gdb14 (\(44{,}00\,\%\)), gdb6
(\(42{,}28\,\%\)) y gdb5 (\(41{,}11\,\%\)). El caso de gdb10, con un gap de
\(52\,\%\) (costo SA \(= 418\) frente a BKS \(= 275\)), es el más severo. Una
hipótesis plausible es que instancias como gdb10, gdb5 y gdb6 presentan estructuras
topológicas donde la solución óptima requiere distribuciones de rutas no alcanzables
mediante movimientos locales simples desde la solución inicial: los operadores de
vecindad no están generando movimientos que escapen de la cuenca de atracción del
mínimo local inicial. Adicionalmente, en instancias con costos absolutos pequeños
(como gdb14, con BKS \(= 100\)), un error de pocos arcos mal asignados puede
traducirse en gaps relativos grandes, amplificando la apariencia del problema.

% --------------------------------------------------------------------------
\subsection{Conclusión}
\label{subsec:conclusion_final}

El SA, con la parametrización evaluada en este experimento, \textbf{garantiza
factibilidad completa y produce soluciones en tiempos del orden de segundos} para el
conjunto de instancias pequeñas (GDB + KSHS). El gap promedio de \(30{,}54\,\%\)
respecto a los BKS publicados indica que el algoritmo converge establemente a
soluciones admisibles, pero que estas se encuentran a una distancia sustancial de los
óptimos de referencia. La mediana de \(27{,}59\,\%\) sugiere que la mitad de las
instancias se resuelven con un gap por debajo de ese umbral, aunque la cola alta (cinco
instancias por encima del \(40\,\%\)) revela que el SA en su forma actual no es
uniformemente competitivo.

La configuración recomendada a partir de este grid search es \(\boldsymbol{\alpha =
0{,}99}\) \textbf{con} \(\boldsymbol{p_{\text{inter}}}\) \textbf{en el rango
\(0{,}6\)--\(0{,}7\)}: son las celdas con menor gap promedio, y la diferencia entre
ellas es inferior a \(0{,}1\) puntos porcentuales. La robustez de esta recomendación
está limitada por la insensibilidad documentada de \(p_{\text{inter}}\), por lo que la
prioridad de calibración futura debe recaer sobre \(\alpha\) y sobre el presupuesto de
búsqueda.

El comportamiento de la parada temprana (mediana de 67 enfriamientos frente al máximo
de 348) apunta a una causa estructural del gap elevado: el criterio
\texttt{max\_reheats\_sin\_mejora} \(= 5\) termina la búsqueda antes de que el
algoritmo haya agotado su capacidad exploratoria real. Las líneas de mejora más directas
son:

\begin{enumerate}
  \item Incrementar \texttt{max\_reheats\_sin\_mejora} a \(10\)--\(15\) para dar más
        presupuesto a instancias difíciles sin aumentar el tiempo en instancias fáciles.
  \item Explorar valores de \(\alpha\) superiores a \(0{,}99\) (por ejemplo, \(0{,}995\)
        o \(0{,}999\)) para verificar si la tendencia monótona observada continúa.
  \item Añadir una fase de búsqueda local determinista (2-opt exhaustivo o
        Lin--Kernighan adaptado a arcos) al final de la ejecución del SA para pulir la
        mejor solución encontrada.
  \item Combinar el SA con una metaheurística complementaria con mayor capacidad de
        diversificación ---como Búsqueda Tabú o Colonia de Abejas--- en un esquema
        híbrido o de reinicio, aprovechando que el SA ya produce soluciones factibles de
        calidad moderada en tiempos reducidos.
\end{enumerate}
